%% parts/p12/interactive-proof.tex --- Part 12.5
\chapter{Interactive and Mechanized Proof}
\label{sec:p12-interactive-proof}

\noindent\emph{Reader note.}
This chapter is a separately authored editorial survey after the exact
inherited Mechanized Proof Checking chapter.  Its house taxonomy, product
landscape, and landmark tables retain their own labels and OPEN evidence
standing; they neither replace nor alter the source definitions, audit lists,
example, warning, or exercises.

\section{Checking, search, and trust are different layers}

\begin{hbdefinition}{def:p12-mechanization-layers}{Proof object, proof term, proof script, and checking layers}
A \emph{proof object} is finite formal data whose declared checker can verify
against a stated judgment.  A \emph{kernel} is the implementation component
that checks or constructs only judgments licensed by one declared foundation.
An \emph{elaborator} translates surface syntax, implicit arguments, notation,
and constraints into kernel input.  A \emph{tactic} is a search program that
constructs a candidate proof object or invokes trusted primitive constructors.
In a calculus with proof terms, a \emph{proof term} is a term whose typing
judgment certifies the represented proposition.  A \emph{proof script} is a
source-level sequence of commands interpreted by elaboration or tactic
machinery to construct or check a proof object; the script need not itself be
the object accepted by the kernel.  Thus a proof term is formal syntactic data
of the calculus, whereas a proof script is an instruction artifact.
The \emph{trusted base} of a reported result is the exact set of axioms,
definitions, checker components, imported artifacts, external oracles, and
hardware/software assumptions on which acceptance of that result depends.

These names are necessary because the surface script, the search that found
it, the explicit object, the final checker, and the assumptions trusted by a
particular run can differ.  Proof objects are formal data, and proof terms are
the term-shaped subclass licensed by a typing judgment; proof scripts are
instruction artifacts; kernels, elaborators, and tactics are implementation
components with different checking or construction roles; and a trusted base
is a result-relative dependency set.  The distinctions are used in the
landscape, mechanized-result, landmark, and release-audit tables below
\cite{Harrison2009,Pollack1998}.
\index{proof object}\index{kernel}\index{elaborator}\index{tactic}
\index{proof term}\index{proof script}\index{trusted base}
\end{hbdefinition}

\begin{hbdefinition}{def:p12-deep-shallow-embedding}{Deep and shallow embeddings}
A \emph{deep embedding} represents the object language by explicit host-level
syntax and defines its judgments over that representation.  A
\emph{shallow embedding} maps object expressions directly to host-language
operations or propositions.  These names are necessary to distinguish retained
object syntax from reused host semantics; both are representation strategies,
not claims of adequacy.  A deep embedding exposes syntax induction but requires
an interpretation, whereas a shallow embedding inherits host behavior but
requires a separate correspondence argument showing that the intended object
distinctions were preserved.  The distinction is used in the encoding audit.
\index{deep embedding}\index{shallow embedding}
\end{hbdefinition}

\begin{hbexample}{def:p12-deep-shallow-conjunction}{Deep and shallow conjunction obligations}
For a source fragment containing propositional conjunction, the two
representations impose different adequacy obligations:
\begin{center}
\small
\begin{tabular}{@{}p{.16\textwidth}p{.31\textwidth}p{.43\textwidth}@{}}
\toprule
Embedding & Representation & Owed adequacy obligation \\ \midrule
Deep & An explicit formula datatype with a conjunction constructor & For
every source formula and source valuation, the deep evaluator returns true on
that formula's explicit representation if and only if the valuation satisfies
the source formula; the representation must also preserve every binding and
syntactic distinction used by the source fragment. \\
Shallow & A direct map from object formulas to host propositions, taking
object conjunction to host conjunction & For every source formula and source
valuation, the resulting host proposition holds if and only if that valuation
satisfies the source formula; host equality or reduction must not identify
object distinctions required later. \\
\bottomrule
\end{tabular}
\end{center}
These quantified biconditionals are owed correspondence statements, not
adequacy theorems established by the notation itself.  The encoding audit
below records which obligation a mechanization actually discharges.
\end{hbexample}

\begin{hbdefinition}{def:p12-lcf-architecture}{LCF architecture}
An \emph{LCF architecture} is the programming discipline in which values of an
abstract theorem type can be constructed only by a small set of trusted
primitive inference functions, while untrusted tactics compose those
functions.  The identifier \(\mathsf{LCF}\) is necessary to name that
constructor boundary, stands for an implementation architecture rather than a
logic, and is not synonymous with exporting a portable proof object.  The
distinction is used in Theorem~\ref{thm:p12-lcf-boundary} and
Table~\ref{tab:p12-assistant-landscape}
\cite{GordonMilnerWadsworth1979}.
\index{LCF architecture}
\end{hbdefinition}

In Table~\ref{tab:p12-assistant-landscape}, the final column reports whether a
particular result can be reduced to a fully explicit object accepted by a
small checker for the declared formal system, with that checker independently
inspectable or independently implementable.  This is an editorial checklist
field, not a technical criterion attributed to de Bruijn and not a theorem or
property inferred from an LCF product label \cite{DeBruijn1970,Pollack1998}.

\begin{hbtheorem}{thm:p12-lcf-boundary}{cited=GordonMilnerWadsworth1979}{Conditional LCF soundness boundary}
If an implementation preserves the abstraction boundary of its theorem type,
and every trusted primitive constructor realizes a sound rule of the declared
calculus, then every theorem value produced by any composition of those
constructors is derivable in that calculus.  The statement does not prove the
primitive implementation correct, identify the whole trusted base, or imply
that a portable certificate is exported.
\end{hbtheorem}

\begin{hbexample}{def:p12-propositions-as-types-example}{A checked propositions-as-types instance}
In the arrow-only calculus of Section~\ref{sec:normalization-curry-howard-02}, let
\(\Theta_{\mathrm{ty}}\) contain \(f{:}\alpha\to_{\mathrm{ty}}\beta\) and
\(x{:}\alpha\).  The concrete judgment
\[
 \Theta_{\mathrm{ty}}\derives_{\mathsf{STLC}}fx{:}\beta
\]
is checked using \hyperref[sheet:stlc]{\texttt{sheet:stlc}}.  Under the
calculus-relative Curry--Howard correspondence of
Definition~\ref{def:p07-curry-howard}, the term records use of an implication
from \(\alpha\) to \(\beta\).  This is one typed example, not a displayed rule
schema and not a blanket identification of all proofs with all programs.
\end{hbexample}

\section{Pure type systems and inductive foundations}

\begin{hbdefinition}{def:p12-pts-lambda-cube}{Pure type system and lambda cube}
A \emph{pure type system} has Barendregt's specification
\(S=(\mathcal S,\mathcal A,\mathcal R)\), where \(\mathcal S\) is a set of
sorts, \(\mathcal A\subseteq\mathcal S\mathbin{\times}\mathcal S\) is a set
of sort-typing axioms \(s_1{:}s_2\), and
\(\mathcal R\subseteq\mathcal S\mathbin{\times}\mathcal S
\mathbin{\times}\mathcal S\) is a set of triples declaring which
dependent-product sorts may be formed: \((s_1,s_2,s_3)\in\mathcal R\)
licenses a dependent product of sort \(s_3\) when its domain has sort \(s_1\)
and its codomain, under the bound variable, has sort \(s_2\).  Its calculus is
denoted
\(\lambda S=\lambda(\mathcal S,\mathcal A,\mathcal R)\).  The identifier
\(\mathsf{PTS}\) names this parameterized family, not one theorem prover.

For the \emph{lambda cube}, fix
\(\mathcal S=\{\ast,\Box\}\),
\(\mathcal A=\{\ast{:}\Box\}\), and the base product pair
\((\ast,\ast)\), where Barendregt abbreviates
\((s_1,s_2,s_2)\) by \((s_1,s_2)\).  The eight vertices arise by adjoining
any subset of \((\ast,\Box)\), \((\Box,\ast)\), and
\((\Box,\Box)\).  These axes respectively permit types to depend on terms,
terms to depend on types, and types to depend on types.  Their standing as
independent formation choices distinguishes the cube from a linear strength
ordering and is used in Table~\ref{tab:p12-lambda-cube}.  Barendregt's
strongest-corner identifier is \(\lambda\mathrm P\omega=\lambda\mathrm C\),
the Calculus of Constructions; it is not CIC
\cite{Barendregt1992,CoquandHuet1988}.
\index{pure type system}\index{lambda cube}\index{Calculus of Constructions}
\end{hbdefinition}

\begin{longtable}{@{}p{.23\textwidth}p{.32\textwidth}p{.35\textwidth}@{}}
\caption{The eight lambda-cube specifications.}
\label{tab:p12-lambda-cube}\\
\toprule
Source identifier & Enabled axes & Boundary \\ \midrule
\endfirsthead
\toprule
Source identifier & Enabled axes & Boundary \\ \midrule
\endhead
\(\lambda{\to}\) & none & Simply typed base corner \\
\(\lambda\mathrm P\) & types on terms & Dependent types \\
\(\lambda 2\) & terms on types & System-F corner \\
\(\lambda\underline{\omega}\) & types on types & Weak \(\lambda\omega\) corner \\
\(\lambda\mathrm P2\) & types on terms; terms on types & Dependent polymorphism \\
\(\lambda\mathrm P\underline{\omega}\) & types on terms; types on types &
Dependent type operators \\
\(\lambda\omega\) & terms on types; types on types & System \(\mathrm F_\omega\) corner \\
\(\lambda\mathrm P\omega=\lambda\mathrm C\) & types on terms; terms on types;
types on types & Calculus of Constructions \\
\bottomrule
\end{longtable}

\begin{hbtheorem}{thm:p12-lambda-cube-classification}{cited=Barendregt1992}{Lambda-cube classification}
For the two-sort PTS convention fixed in
Definition~\ref{def:p12-pts-lambda-cube}, the independent three formation axes
give exactly the eight vertices in Table~\ref{tab:p12-lambda-cube};
\(\lambda\mathrm C=\lambda\mathrm P\omega\) is the vertex permitting all
three.
\end{hbtheorem}

\begin{hbdefinition}{def:p12-inductive-type-boundaries}{Universes, inductives, W-types, and two equalities}
A \emph{universe} is a type whose elements code types admitted at a declared
level; a hierarchy and its cumulativity relation must be stated rather than
inferred from the word.  An \emph{inductive type} is introduced by a
constructor specification admitted by the formal system's own explicitly
stated inductive-definition or positivity criterion, together with that
system's associated elimination principles; a report must name the actual
criterion rather than importing a generic one.  A \emph{W-type} is the
well-founded-tree type
generated by a family of node labels and arities; an assistant may encode such
trees through a general inductive facility instead of exposing W-types as a
primitive.  \emph{Definitional equality} is the kernel conversion relation
used during checking, whereas \emph{propositional equality} is represented by
an inhabitant of a declared identity/equality type.

These names are needed to prevent broad dependent-type-theory family labels
and CIC from hiding
different universe, positivity, elimination, and equality choices.  They stand
for type formers or judgmental/propositional relations, not for interchangeable
marketing categories, and the distinctions control the foundation cells of
Table~\ref{tab:p12-assistant-landscape}.
\index{universe}\index{inductive type}\index{W-type}
\index{definitional equality}\index{propositional equality}
\end{hbdefinition}

\begin{hbdefinition}{def:p12-cic-boundary}{CIC and pCIC survey boundary}
The \emph{Calculus of Inductive Constructions} (\(\mathsf{CIC}\)) denotes
CoC extended by a precisely declared universe/cumulativity discipline,
strictly controlled inductive definitions, and their elimination principles;
\emph{pCIC} additionally records the polymorphic cumulative presentation used
by current Rocq documentation.  The identifiers are necessary to distinguish
these extensions from the CoC cube corner, stand for source-defined calculus
families rather than every dependent type theory, and are not names for Lean's
exact kernel theory.  They are used only to classify the landscape; formal
calculus surfaces remain in Part~7 or in version-specific manuals
\cite{PaulinMohring1993,PaulinMohring2015,RocqManual2026}.
\index{Calculus of Inductive Constructions}\index{pCIC}
\end{hbdefinition}

No blanket normalization, decidability, or consistency theorem is asserted
for dependent-type-theory families or CIC as a whole.  Those conclusions
depend on the exact constructors, universe rules, conversion relation, and
enabled extensions.

\begin{hbdefinition}{def:p12-computational-reflection}{Computational reflection in mechanized proof}
A mechanized proof step uses \emph{computational reflection} for a declared
formula fragment when it represents formulas of that fragment inside the
prover, runs a procedure on those representations, and applies a previously
kernel-checked soundness theorem connecting a successful computation to the
represented proposition.  The term therefore names a proof method, not a
primitive logical rule, a synonym for tactic execution, or permission to trust
an unchecked oracle.  Its nonredundant feature is the reusable internal
soundness theorem together with a checked evaluation path.  This distinction
is required by the mechanized-result report, which records the representation
and interpretation, the soundness theorem, and whether evaluation is kernel
reduction, proof-producing code, or an externally trusted path
\cite{Harrison1995Reflection}.
\index{computational reflection}\index{reflection!computational}
\end{hbdefinition}

\section{Proof-assistant landscape}

\begin{landscape}
\thispagestyle{plain}
\begingroup\scriptsize
\setlength{\tabcolsep}{3pt}
\begin{longtable}{@{}p{1.00in}p{1.85in}p{2.10in}p{2.15in}@{}}
\caption{Proof-assistant and verifier landscape.  The final
column records the editorial small-checker question stated above; it is not a
product-quality score or an attributed source criterion.}
\label{tab:p12-assistant-landscape}\\
\toprule
Tool/version & Declared foundation & Checking and automation/trust path &
Independent small-checker path \\ \midrule
\endfirsthead
\toprule
Tool/version & Declared foundation & Checking and automation/trust path &
Independent small-checker path \\ \midrule
\endhead
Lean 4.32.1 & Dependent type theory with universes, inductives, quotients, and
documented proof irrelevance features & Elaborator and tactics produce terms
for the kernel; imported axioms, code-backed decisions, and external artifacts
must be reported & Qualified for ordinary kernel-checked terms; the exact axiom
surface and any native/oracle path remain result-specific
\cite{LeanManual2026,DeMouraUllrich2021} \\ \addlinespace
Rocq 9.1.1 & Polymorphic cumulative calculus of inductive constructions &
Tactics elaborate proof terms checked by the kernel; plugins, axioms, and
extraction are separate trust questions & Qualified for exported kernel terms
relative to the declared environment; extraction is not itself a proof
certificate \cite{RocqManual2026,PaulinMohring2015} \\ \addlinespace
Isabelle2025-2 & Generic logical framework; Isabelle/HOL instantiates it with
classical higher-order logic as an object logic & LCF-family trusted inference kernel with proof
methods and external automation reconstructed or otherwise reported &
Session theorem values have an LCF trust path; a portable independently
checkable object depends on the actually used export/reconstruction path
\cite{Isabelle2025} \\ \addlinespace
Agda 2.8.0 & Baseline dependent type theory with
\texttt{--no-sized-types}; neither \texttt{--cubical} nor
\texttt{--rewriting} is enabled in this row & Elaborator/type checker verifies
definitions; reflection and enabled language options belong in the report &
Configuration-dependent: this baseline names its disabled optional extensions;
any other result must record its exact options and postulates before claiming
an independent small-checker path \cite{AgdaReference2026} \\ \addlinespace
HOL Light & Classical simple type theory & Small OCaml LCF kernel; derived
rules and automation construct protected theorem values & LCF discipline is
explicit, but portability of a complete proof object is a separate export
claim \cite{HOLLightSite2026} \\ \addlinespace
HOL4 Trindemossen-2 & Classical higher-order logic & LCF-family kernel and
derived proof tools; external components must be reported & LCF trust path is
documented; independent portable-object status depends on the chosen workflow
\cite{HOL4Site2026} \\ \addlinespace
PVS 7.1 & Higher-order logic with dependent predicate subtypes & Type checking,
proof commands, decision procedures, and integrated automation cooperate &
Not established merely by the product description; report any exported proof
and its independent checker \cite{PVS71Documentation} \\ \addlinespace
Mizar 8.1.15 / MML 5.94.1493 & First-order set theory over
Tarski--Grothendieck set theory & Verifier checks elaborated articles relative
to the declared MML environment & Relative to the verifier and exact imported
environment; no LCF theorem datatype is inferred
\cite{MizarLibrary2026} \\ \addlinespace
ACL2 8.7 & First-order logic of total recursive functions in an executable
applicative Lisp setting & Integrated prover uses rewriting, induction, and
verified or trusted extensions recorded by the event history & An exported
small-checker proof object is not established by the cited portal; report the
actual certification path \cite{ACL2Documentation2026} \\ \addlinespace
Metamath, retrieved 2026-08-11 & The database itself declares constants,
variables, axioms, and rules & A small verifier checks explicit substitutions
and proof steps; the system is not intrinsically tied to ZFC & Meets the
operational pattern relative to the exact database and verifier, whose hash or
version must accompany the result \cite{MetamathSite2026} \\
\bottomrule
\end{longtable}
\endgroup
\end{landscape}

\section{Landmark formalizations}

\begin{landscape}
\thispagestyle{plain}
\begingroup\scriptsize
\setlength{\tabcolsep}{3pt}
\begin{longtable}{@{}p{1.05in}p{1.20in}p{2.10in}p{2.75in}@{}}
\caption{Landmark formalizations and exact claim boundaries.}
\label{tab:p12-landmark-formalizations}\\
\toprule
Result & Assistant/foundation & Formal result & Trust and scope boundary \\
\midrule
\endfirsthead
\toprule
Result & Assistant/foundation & Formal result & Trust and scope boundary \\
\midrule
\endhead
Four-color theorem & Coq/Rocq, CIC-family development & A machine-checked
formalization of the four-color theorem & Does not claim that every historical
paper proof was checked; imports and the kernel version remain part of the
report \cite{Gonthier2008} \\ \addlinespace
Odd-order theorem & Coq/Rocq, constructive type theory & A machine-checked
formalization of the Feit--Thompson odd-order theorem & Claim is relative to the
declared development and libraries, not to all group-theory presentations
\cite{GonthierEtAl2013} \\ \addlinespace
Flyspeck & HOL Light and Isabelle components, classical HOL & A formal proof of
the Kepler conjecture & Both assistants and the correspondence between checked
components must be reported; no single product label erases that boundary
\cite{HalesEtAl2017} \\ \addlinespace
CompCert & Coq/Rocq, CIC-family development & Semantic-preservation theorems
for the stated compiler passes, source languages, and targets & Does not assert
that every C program, compiler option, assembler, or hardware behavior is
bug-free \cite{Leroy2009} \\ \addlinespace
seL4 & Isabelle/HOL & Functional correctness from the stated abstract
specification to the verified C implementation & Relies on the cited compiler,
assembler, hardware, and model assumptions and does not establish every
security or timing property \cite{KleinEtAl2009} \\
\bottomrule
\end{longtable}
\endgroup
\end{landscape}

For the historical placement of these developments, see the completed
chronology and capsules in Chapter~\ref{sec:p13-chronology-and-biography}.

\begin{proofboundary}{def:p12-hott-pointer}{HoTT and univalence pointer}
The formal surface meant here is intensional Martin--L\"of type theory with
dependent products, dependent sums, identity types, and a cumulative universe
discipline, extended in the cited presentation by the univalence axiom and
higher inductive types \cite{HoTTBook2013}. Its exact type-forming and equality
rules are not interchangeable with the narrower simply typed or System-F
sheets. The calculus and notation boundary is recorded in
Definition~\ref{def:p07-mltt-boundary} and
\ref{def:p07-hott-cubical-boundary}; this pointer asserts no normalization or
canonicity theorem for the extension.
\end{proofboundary}

\begin{hbdefinition}{def:p12-program-extraction}{Program extraction boundary}
For the Coq-style pipeline considered here, the source is a well-typed term of
the Calculus of Inductive Constructions. Extraction erases terms whose types
live in proof-irrelevant \(\mathsf{Prop}\), removes type information not needed
at run time, and translates the remaining computational term to a declared
functional target such as OCaml. The name is necessary to distinguish erasure
and translation from synthesis of a new algorithm. It stands as a compiler
pass with a source/target relation, differs because proof components may be
removed only under the calculus's elimination restrictions, and is used in the
correspondence and trust record \cite{Letouzey2008}.
\index{program extraction}\index{proof erasure}
\end{hbdefinition}

The correspondence claim is conditional: evaluation of the extracted target
must simulate the computational content of the well-typed source under the
specified erasure theorem. A complete evidence record names the kernel and
library assumptions used to type-check the source, extraction implementation
and version, target compiler, runtime, foreign primitives, and the exact input
and output. Kernel checking of the source theorem does not verify an unproved
compiler optimization or external runtime.

\begin{hbdefinition}{def:p12-hammer-pipeline}{Hammer and proof reconstruction}
A proof-assistant hammer performs four distinct stages: premise selection from
the local library; translation of the goal and selected facts to an external
logic; search by one or more automated provers; and reconstruction or replay
of returned evidence in the assistant's kernel. The name is needed because an
external prover's success token and a kernel theorem have different trust
standing. It denotes a metalinguistic automation pipeline, differs from an
internal tactic by crossing a translation/trust boundary, and is used to state
the acceptance rule: only successful kernel reconstruction closes the proof
\cite{BlanchetteEtAl2016Hammering}.
\index{hammer!proof assistant}\index{proof reconstruction}\index{premise selection}
\end{hbdefinition}

When reconstruction fails, the external search result remains diagnostic
evidence rather than a theorem. A reproducible hammer record therefore keeps
the translated problem, selected premise identities, prover and limits,
returned proof or core, reconstruction command, assistant version, imports,
and final kernel result.
